\begin{description}
\item[a.)] Intensity
  \[ I = \frac{L}{4\pi D^2} \]
  \hfill(1 point)
\item[b.)] Absorption rate
  \begin{align*}
  \mathcal{A} &= \left(1-\alpha\right)\pi R^2 I\\
              &= \left(1-\alpha\right)\frac{LR^2}{4D^2}
  \end{align*}
  \hfill(1 point)
\item[c.)] Light energy reflected by the planet per unit time is
  \[ \alpha\pi R^2 I = \frac{\alpha L R^2}{4D^2} \]
  \hfill(1 point)

  Hence the planet's luminosity is
  \[ L_{\mathrm{planet}} = \frac{\alpha L R^2}{4D^2} \]
  \hfill(1 point)
\item[d.)] Here, we will neglect the planet's internal source of energy. Let
  $T$ be the black-body temperature of the planet's surface in kelvins. Since
  the planet is rotating fast, we may assume that its surface is being heated
  up uniformly to approximately the same temperature $T$. The total amount of
  black-body radiation emitted by the planet's surface is from
  Stefan-Boltzmann law: $4\pi R^2 \cdot \sigma T^4$, $\sigma$ being
  Stefan-Boltzmann constant. At equilibrium, that is when the temperature
  remains steady, this emission rate must be equal to the absorption rate in
  b.). \hfill(1 point)

  Hence
  \begin{align*}
  4\pi R^2 \cdot \sigma T^4 &= \left(1-\alpha\right)\frac{LR^2}{4D^2}\\
  T &= \left[\left(1-\alpha\right)\frac{L}{16\pi\sigma D^2}\right]^{\frac14}
  \end{align*}
  \hfill(1 point)
\item[e.)] In this case the emitted black-body radiation is mostly from the
  planet's surface facing the star. The emitting surface area is now only
  $2\pi R^2$ and not $4\pi R^2$. Hence the surface temperature is given by
  $T'$, where
  \[ 2\pi R^2 \cdot \sigma \left(T'\right)^4
     = \left(1-\alpha\right)\frac{LR^2}{4D^2} \]
  \hfill(1 point)
  \[ T' = \left[\left(1-\alpha\right)\frac{L}{8\pi\sigma D^2}\right]^{\frac14}
        = \left(2\right)^{\frac14}\cdot T \approx \left(1.19\right)T \]
  \hfill(1 point)
\item[f.)]
  \begin{align*}
  T &= \left[\left(1-\alpha\right)\frac{L}{16\pi\sigma D^2}\right]^{\frac14}\\
  T &= \left[\left(1-0.25\right)\times
       \frac{3.826\times 10^{26}}
            {16\pi\times 5.67\times 10^{-8}\times
             \left(1.523\times 1.496\times 10^{11}\right)^2}\right]^{\frac14}\\
    &= 209.8 \simeq 210\ \mathrm{K} = -63\,^\circ\mathrm{C}
  \end{align*}
  \hfill(2 points)
\end{description}
